\documentclass{beamer}

\usepackage{graphicx}
\usepackage{textpos}
\usepackage{listings}
\usepackage{lstautogobble}

\usetheme{Madrid}
\useoutertheme{miniframes} % Alternatively: miniframes, infolines, split

% Setup the university's color pallette
\definecolor{UIUCorange}{RGB}{19, 41, 75} % UBC Blue (primary)
\definecolor{UIUCblue}{RGB}{232, 74, 39} % UBC Grey (secondary)

\definecolor{codegreen}{rgb}{0,0.6,0}
\definecolor{codegray}{rgb}{0.5,0.5,0.5}
\definecolor{codepurple}{rgb}{0.58,0,0.82}
\definecolor{backcolour}{rgb}{0.95,0.95,0.92}

\lstdefinestyle{python}{
    backgroundcolor=\color{backcolour},   
    commentstyle=\color{codegreen},
    keywordstyle=\color{magenta},
    numberstyle=\tiny\color{codegray},
    stringstyle=\color{codepurple},
    basicstyle=\ttfamily\footnotesize,
    breakatwhitespace=false,         
    belowskip=-0.5em,
    breaklines=true,                 
    captionpos=b,                    
    keepspaces=true,                 
    %numbers=left,                    
    numbersep=5pt,                  
    showspaces=false,                
    showstringspaces=false,
    showtabs=false,                  
    tabsize=2
}

\lstset{style=python}

\AtBeginSection[]{
    \begin{frame}
        \vfill
        \centering
        \begin{beamercolorbox}[sep=8pt,center,shadow=true,rounded=true]{title}
            \usebeamerfont{title}\insertsectionhead\par%
        \end{beamercolorbox}
        \vfill
    \end{frame}
}
% Setup the university's color pallette
\definecolor{UIUCorange}{RGB}{19, 41, 75} % UBC Blue (primary)
\definecolor{UIUCblue}{RGB}{232, 74, 39} % UBC Grey (secondary)


\setbeamercolor{palette primary}{bg=UIUCorange,fg=white}
\setbeamercolor{palette secondary}{bg=UIUCblue,fg=white}
\setbeamercolor{palette tertiary}{bg=UIUCblue,fg=white}
\setbeamercolor{palette quaternary}{bg=UIUCblue,fg=white}
\setbeamercolor{structure}{fg=UIUCorange} % itemize, enumerate, etc
\setbeamercolor{section in toc}{fg=UIUCblue} % TOC sections

\setbeamercolor{subsection in head/foot}{bg=UIUCorange,fg=UIUCblue}
\setbeamercolor{subsection in head/foot}{bg=UIUCorange,fg=UIUCblue}

\usepackage[utf8]{inputenc}
\usepackage{graphicx}

%Information to be included in the title page:
\title{\textbf{Classes}}
\author{\textbf{David H Smith IV}}
\institute[\textbf{UIUC}]{\textbf{University of Illinois Urbana-Champaign}}
\date{\textbf{Tues, Nov 16 2021}}

\setbeamertemplate{title page}[default][colsep=-4bp,rounded=true]
\addtobeamertemplate{title page}{\vspace{3\baselineskip}}{}
\addtobeamertemplate{title page}{
    \begin{textblock*}{\paperwidth}(-1.0em, -1.2em)
        \includegraphics[width=\paperwidth, height=\paperheight]{imgs/uiuc.png}
    \end{textblock*} 
}{}

\begin{document}

\frame{\titlepage}

\section{Reminders}

%
% Slide 1
%
\begin{frame}
    \frametitle{Reminders}
    \begin{itemize}
        \item 
    \end{itemize}
\end{frame}

\section{The Birds Eye View}

%
% Slide 2
%
\begin{frame}[fragile]
    \frametitle{Classes, Objects, Instances, oh my!}
    \begin{enumerate}[A]
        \item \textbf{Classes:} The actual Python code that provides instructions on how to build a class (\lstinline|__init__()|), the attributes in the class, and definitions for the class functions.
            \pause
            \begin{itemize}
                \item \textbf{Class Attribute: } A value in the class that is accessible to all instances of that class.
                    \pause
                \item \textbf{Instance Attribute: } A value that is only accessible to a given instance.
                    \pause
                \item \textbf{Instance Method: } A function that is callable from within a given instance. The words `method' and `function' mean the same thing.
            \end{itemize}
            \pause
        \item \textbf{Instance:} An object that was created using a given class. We can have multiple instances of the same class.
            \pause
        \item \textbf{Object:} Just another name for the thing we instantiated.
    \end{enumerate}
\end{frame}

\section{Making a Class}

%
% Slide 2
%
\begin{frame}[fragile]
    \frametitle{Making a Class: Poll Question}
    Which of the following is the correct way of instantiating the class \lstinline|Foo|?
    \begin{lstlisting}[language=Python, autogobble]
    class Foo:
        def __init__(self):
            print("I'm a class!")
    \end{lstlisting}
    \vfill
    \begin{enumerate}[A]
        \item \lstinline|foo = Foo()|
        \item \lstinline|foo = Foo(self)|
        \item \lstinline|foo = Foo.__init__()|
        \item \lstinline|foo = __init__()|
    \end{enumerate}
\end{frame}

%
% Slide 2
%
\begin{frame}[fragile]
    \frametitle{Making a Class: Poll Question}
    What is the result of the following code?
    \begin{lstlisting}[language=Python, autogobble]
    class Foo:
        def bar(self, x, y):
            return x + y

    foo = Foo()
    x = foo.bar(1, 2)
    print(x)
    \end{lstlisting}
    \vfill
    \begin{enumerate}[A]
        \item SyntaxError
        \item NameError
        \item 3
        \item TypeError
    \end{enumerate}
\end{frame}


%
% Slide 2
%
\begin{frame}[fragile]
    \frametitle{Making a Class: \lstinline|self|}
    What is the result of the following code?
    \begin{lstlisting}[language=Python, autogobble]
    class Foo:

        def __init__(self):
            print("I'm a class!")

        def get_id(self):
            return id(self)

    foo = Foo()
    print(id(foo) == foo.get_id())
    \end{lstlisting}
    \vfill
    \begin{enumerate}[A]
        \item True
        \item False
        \item Trick question
    \end{enumerate}
\end{frame}

%
% slide 2
%
\begin{frame}[fragile]
    \frametitle{\lstinline|self| as an automatic first argument}
    Given this class...\\
    \vfill
    \begin{lstlisting}[language=python, autogobble]
    class Foo:
        def __init__(self):
            print("class created!")

        def bar(self, x, y):
            return x + y
    \end{lstlisting}
    \pause
    \vspace{0.5cm}
    This...
    \vfill
    \begin{lstlisting}[language=python, autogobble]
    f = Foo()
    x = f.bar(5, 6)
    \end{lstlisting}
    \pause
    \vspace{0.5cm}
    Is the same as this...\\
    \vfill
    \begin{lstlisting}[language=python, autogobble]
    f= Foo()
    x = Foo.get_id(f, 5, 6)
    \end{lstlisting}
\end{frame}

%
% slide 2
%
\begin{frame}[fragile]
    \frametitle{Classses in General}
    You've seen this before, you just didn't know it... Any thoughts?
    \pause
    \begin{lstlisting}[language=python, autogobble]
    lst = list()
    s = set()
    d1 = dict()
    \end{lstlisting}
    \vfill
    \begin{minipage}{0.49\textwidth}
        \pause
        These...
        \begin{enumerate}
            \item \lstinline|lst.append(x)|
            \item \lstinline|s.add(x)|
            \item \lstinline|d1.update(d2)|
        \end{enumerate}
    \end{minipage}
    \hfill
    \begin{minipage}{0.49\textwidth}
        \pause
        Are the same as these
        \begin{enumerate}
            \item \lstinline|list.append(lst, x)|
            \item \lstinline|set.add(s, x)|
            \item \lstinline|dict.update(d1, d2)|
        \end{enumerate}
    \end{minipage}
    \vfill
    \pause
    \textbf{This is more syntactic sugar brought to you by Python. The key takeaway is that \lstinline|self| is automatically added in front and refers to the object before the dot (\lstinline|.|).}
\end{frame}


%
% slide 2
%
\begin{frame}[fragile]
    \frametitle{Making Classes: Key Takeaways}

    \begin{enumerate}
        \item Classes are a list of instructions for how to instantiate an object just as functions are a list of instructions on how to perform an operation given some data.
            \pause
        \item Classes are abstract descriptions, objects are concrete and actually exist.
            \pause
        \item \lstinline|__init__| Is called when you create a function but is never explicitly called.
            \begin{itemize}
                \item Example: \lstinline|foo = Foo()|
            \end{itemize}
            \pause
        \item Self is automatically passed in and refers to the object bound to the variable before the dot.
            \begin{itemize}
                \item Example: \lstinline|foo.call_function()|
            \end{itemize}
            \pause
        \item \lstinline|__init__| is not required. If it is not present in a class definition a default one will be provided and used to instantiate the object.
    \end{enumerate}

\end{frame}

\section{Attributes and Methods}

%
%
%
\begin{frame}[fragile]
    \frametitle{Poll Question: }
    \begin{minipage}{0.35\textwidth}
        \begin{lstlisting}[language=Python, autogobble, basicstyle=\tiny]
        class Name:
            name_count = 0
            def __init__(self, name):
                ??
                self.name = name

        n1 = Name("foo")
        n2 = Name("bar")
        n3 = Name("baz")
        \end{lstlisting}
    \end{minipage}
    \hfill
    \begin{minipage}{0.59\textwidth}
        Which of the following lines can be used to increment the class attribute count of Name instances that have been instantiated?
        \begin{itemize}
            \item \lstinline|Name.name_count += 1|
            \item \lstinline|self.name_count += 1|
            \item \lstinline|name_count += 1|
            \item \lstinline|self.name_count = Name.name_count + 1|
            \item \lstinline|Name.name_count = self.name_count + 1|
        \end{itemize}
    \end{minipage}
    \vfill
    \textbf{Consider each and then I'll go through and ask true (A) or false (B) for each.}
\end{frame}

%
%
%
\begin{frame}[fragile]
    \frametitle{Poll Question: }
    \begin{minipage}{0.49\textwidth}
        \begin{lstlisting}[language=Python, autogobble]
        class Name:
            name_count = 0
            def __init__(self, name):
                self.name_count += 1
                self.name = name

        n1 = Name("foo")
        n2 = Name("bar")
        n3 = Name("baz")
        print(n1.name_count, n2.name_count, n3.name_count)
        \end{lstlisting}
    \end{minipage}
    \hfill
    \begin{minipage}{0.49\textwidth}
        What is the output of the program on the right?
        \begin{enumerate}[A]
            \item \lstinline|1 1 1|
            \item \lstinline|1 2 3|
            \item \lstinline|3 3 3|
            \item NameError
        \end{enumerate}
    \end{minipage}
\end{frame}

%
%
%
\begin{frame}[fragile]
    \frametitle{Scoping in Python Sucks (I Hate it Very Much)}
    \begin{minipage}{0.47\textwidth}
        This...
        \begin{lstlisting}[language=Python, autogobble, basicstyle=\tiny]
        class Name:
            name_count = 0
            def __init__(self, name):
                self.name_count += 1
                self.name = name
        \end{lstlisting}
    \end{minipage}
    \hfill
    \begin{minipage}{0.52\textwidth}
        Is the same as this.
        \begin{lstlisting}[language=Python, autogobble, basicstyle=\tiny]
        class Name:
            name_count = 0
            def __init__(self, name):
                self.name_count = self.name_count + 1
                self.name = name
        \end{lstlisting}
    \end{minipage}
    \vfill
    \pause
    Looking at the one on the right:
    \begin{enumerate}
        \item Python starts by evaluating the expression on the right.
            \pause
        \item \lstinline|self.name_count + 1|:  \lstinline|self.name_count| isn't an instance attribute so it resolves to the class attribute.
            \pause
        \item \lstinline|self.name_count = 1|: \lstinline|self.name_count| doesn't exist as an instance level attribute so scope resolution decides to create a new instance attribute, thus leaving the class attribute unaffected.
    \end{enumerate}
    \pause
    \textbf{Always use the class name to change class attributes. Bad confusing things happen otherwise.}
\end{frame}

%
%
%
\begin{frame}[fragile]
    \frametitle{Poll Question: }
    \begin{minipage}{0.49\textwidth}
        \begin{lstlisting}[language=Python, autogobble]
        class Name:
            name_count = 0
            def __init__(self, name):
                Name.name_count += 1
                self.name = name

        n1 = Name("foo")
        n2 = Name("bar")
        n3 = Name("baz")
        print(??)
        \end{lstlisting}
    \end{minipage}
    \hfill
    \begin{minipage}{0.49\textwidth}
        Which of the following lines CANNOT be used to identify how many name\_count were created?
        \begin{enumerate}[A]
            \item \lstinline|Name.name_count|
            \item Either \lstinline|n1.name_count| or \lstinline|n2.name_count| or \lstinline|n3.name_count|
            \item \lstinline|Name().name_count|
        \end{enumerate}
    \end{minipage}
    \vfill
    \textbf{Which should be used?}
\end{frame}


%
%
%
\begin{frame}[fragile]
    \frametitle{Poll Question: The Race Class}
    \begin{minipage}{0.69\textwidth}
        \begin{lstlisting}[language=Python, autogobble, basicstyle=\tiny]
        class Racer:

            finished_list = []

            def __init__(self, name, number):
                self.name = name
                self.number = number

            def finished(self):
                Racer.finished_list.append(self)
                print("finished in", len(Racer.finished_list))

        \end{lstlisting}
        \pause
        \vspace{0.5cm}
        What is produced by the following code?
        \vfill
        \begin{lstlisting}[language=Python, autogobble, basicstyle=\tiny]

        r1 = Racer("David", 13)
        r2 = Racer("Dipti", 142)
        print(Racer.finished_list)
        r2.finished()
        r1.finished()
        print([r.name for r in Racer.finished_list])
        \end{lstlisting}
    \end{minipage}
    \hfill
    \begin{minipage}{0.29\textwidth}
        \begin{enumerate}[A]
            \item 
                \begin{lstlisting}[autogobble, basicstyle=\tiny]
                []
                finished in 2
                finished in 1
                ['David', 'Dipti']
                \end{lstlisting}
            \item AttributeError
            \item 
                \begin{lstlisting}[autogobble, basicstyle=\tiny]
                []
                finished in 1
                finished in 2
                ['Dipti', 'David']
                \end{lstlisting}
            \item 
                \begin{lstlisting}[autogobble, basicstyle=\tiny]
                []
                finished in 1
                finished in 1
                []
                \end{lstlisting}
        \end{enumerate}
    \end{minipage}
\end{frame}

%
%
%
\begin{frame}[fragile]
    \frametitle{Poll Question: Ready-to-Go}
    \begin{minipage}{0.69\textwidth}
        \begin{lstlisting}[language=Python, autogobble, basicstyle=\tiny]
        class ReadyToGo:

                ready = 0
                instances = 0

                def __init__(self, name):
                        self.name = name
                        self.ready = False
                        ReadyToGo.instances += 1

                def set_ready(self):
                        ReadyToGo.ready += 1
                        self.ready = True
        \end{lstlisting}
        \pause
        \vspace{0.1cm}
        What is produced by the following code?
        \vfill
        \begin{lstlisting}[language=Python, autogobble, basicstyle=\tiny]
        p1 = ReadyToGo("Alice")
        p2 = ReadyToGo("Bob")
        p3 = ReadyToGo("Charlie")
        players = [p1, p2, p3]

        p1.set_ready()
        p3.set_ready()

        for player in players:
                if not player.ready:
                        print(player.name, "is not ready")
        \end{lstlisting}
    \end{minipage}
    \hfill
    \begin{minipage}{0.29\textwidth}
        \begin{enumerate}[A]
            \item 
                \begin{lstlisting}[autogobble, basicstyle=\tiny]
                Bob is not ready
                \end{lstlisting}
            \item SyntaxError
            \item NameError
            \item AttributeError
            \item 
                \begin{lstlisting}[autogobble, basicstyle=\tiny]
                Alice is not ready
                Bob is not ready
                Charlie is not ready
                \end{lstlisting}
        \end{enumerate}
    \end{minipage}
\end{frame}

\section{Modifying Attributes after Instantiation}
%
%
%
\begin{frame}[fragile]
    \frametitle{Poll Question: }
    \begin{minipage}{0.69\textwidth}
        What is produced by the following code?
        \begin{lstlisting}[language=Python, autogobble, basicstyle=\tiny]
        class Name:
            name_count = 0
            def __init__(self, name):
                self.name = name
                Name.name_count += 1

        foo = Name("foo")
        bar = Name("bar")
        foo.name = "Fred"
        print(bar.name, foo.name)
        \end{lstlisting}
    \end{minipage}
    \hfill
    \begin{minipage}{0.29\textwidth}
        \begin{enumerate}[A]
            \item AttributeError
            \item NameError
            \item \lstinline|bar foo|
            \item \lstinline|bar Fred|
        \end{enumerate}
    \end{minipage}
\end{frame}

%
%
%
\begin{frame}[fragile]
    \frametitle{Poll Question: }
    \begin{minipage}{0.69\textwidth}
        What is produced by the following code?
        \begin{lstlisting}[language=Python, autogobble, basicstyle=\tiny]
        class Name:
            name_count = 0
            def __init__(self, name):
                self.name = name
                Name.name_count += 1

        foo = Name("foo")
        bar = Name("bar")
        Name.name_count = 1000
        print(bar.name_count)
        \end{lstlisting}
    \end{minipage}
    \hfill
    \begin{minipage}{0.29\textwidth}
        \begin{enumerate}[A]
            \item AttributeError
            \item NameError
            \item 1000
            \item 2
        \end{enumerate}
    \end{minipage}
\end{frame}

%
%
%
\begin{frame}[fragile]
    \frametitle{Poll Question: }
    \begin{minipage}{0.69\textwidth}
        What is produced by the following code?
        \begin{lstlisting}[language=Python, autogobble, basicstyle=\tiny]
        class Name:
            name_count = 0
            def __init__(self, name):
                self.name = name
                Name.name_count += 1

        foo = Name("foo")
        bar = Name("bar")
        foo.name_count = 1000
        print(bar.name_count)
        \end{lstlisting}
    \end{minipage}
    \hfill
    \begin{minipage}{0.29\textwidth}
        \begin{enumerate}[A]
            \item AttributeError
            \item NameError
            \item 1000
            \item 2
        \end{enumerate}
    \end{minipage}
\end{frame}

%
% slide 2
%
\begin{frame}[fragile]
    \frametitle{Key Takeaways}
    \begin{enumerate}
        \item How to reference attributes best practices:
            \begin{itemize}
                \item Reference class attributes using the class name:
                    \vspace{-0.2cm}
                    \begin{lstlisting}[language=python, autogobble, basicstyle=\tiny]
                class Foo:
                    count = 0
                    def __init__(self):
                        Foo.count += 1
                    \end{lstlisting}
                    \vfill
                \item Reference instance attributes and methods using self \textit{when inside the class}.
                    \vspace{-0.2cm}
                    \begin{lstlisting}[language=python, autogobble, basicstyle=\tiny]
                class Foo:
                    def __init__(self, name):
                        self.name = name
                    def print_name(self):
                        print(self.name)
                    \end{lstlisting}
                    \vfill
                \item Reference instance attributes and methods using instance's variable name \textit{when outside the class}.
                    \vspace{-0.2cm}
                    \begin{lstlisting}[language=python, autogobble, basicstyle=\tiny]
                class Foo:
                    def __init__(self, name):
                        self.name = name
                x = Foo("bar")
                print(x.name)
                    \end{lstlisting}
                    \vfill
            \end{itemize}
        \item \lstinline|self| must be the first parameter in every instance methods arguments.
    \end{enumerate}
\end{frame}

\end{document}
